% Folder 136, batch 15 (archivists' pages 281-300 as served by the batch PDF).
% Pass: Opus 5.5 (claude-opus-5-5), 2026-10-03 — first pass, unchecked against the pages by a human.
%
% Page numbers follow the batch PDF (page 1 of the batch = 281), so that the
% facsimile turns with the text. The pencilled numbers visible on the leaves
% run about one to two ahead (282 on the first sheet, 302 on the last).
% Skipped as unrelated versos: 281, 283, 285, 287, 289, 291, 293, 295, 297,
% 299 (typescript pages and exercise sheets with no writing of his on the
% mathematics in hand). Page 300 is a cover in his hand naming the next run.

\documentclass[11pt,a4paper]{article}
\input{../preamble/grothendieck}

\folder{136}
\batch{15}
\pages{281}{300}
\foldertitle{Complexe de De Rham à puissance divisée [conférence de 1976 à l’IHÉS] : notes manuscrites (s.d.).}
\dating{[à partir de 1975-1976]}
\watermark{Édition de démonstration}

\begin{document}

\section{Coalgèbres : noyau, intersection des $P_i$, produit tensoriel au-dessus de $S$}

\page{282}
\note{Page de calculs sans prose suivie ; l'argument commence avant ce lot (la ligne du haut s'ouvre sur un trait, comme la suite d'une ligne précédente). Les objets $C$, $C'$, $C''$, $S$ et le foncteur $V$ ne sont pas définis sur ces pages.}

\[
\cdots \to P \longrightarrow C\otimes C' \rightrightarrows C\otimes C'\otimes S
\]

\begin{tikzcd}
V(P) & V(C\otimes C') \arrow[l] \arrow[dl] & V(C)\times V(C')\times V(S) \arrow[l, bend left=15] \arrow[l, bend right=15] \\
V(C'') & &
\end{tikzcd}

\note{dans $V(C\otimes C')$, le signe $\otimes$ est écrit sur un symbole biffé illisible.}

\[
P = \operatorname{Ker}\bigl(C\otimes C' \rightrightarrows C\otimes C'\otimes S\bigr)
\]
\note{le dernier facteur $\otimes S$ est biffé dans cette formule.}

\[
P\otimes S = \operatorname{Ker}\bigl(C\otimes C'\otimes S \rightrightarrows C\otimes S\otimes C'\otimes S\bigr)
\]
\note{précédé de « $P$ \struck{$\otimes$} $\longrightarrow$ » ; des accolades soulignent $C\otimes$ dans le premier terme et $C\otimes S\otimes$ dans le second.}

\[
P \to C\otimes C', \qquad P\otimes S \to C\otimes C'\otimes S
\]

\begin{tikzcd}
P \arrow[r, hook, "i"] \arrow[d, "?"'] & C\otimes C' \arrow[d, bend right=20, "\alpha"'] \arrow[d, bend left=20, "\beta"] \\
P\otimes S \arrow[r, hook, "{i\otimes\mathrm{id}_S}"] & C\otimes C'\otimes S
\end{tikzcd}

\note{la flèche du bas est tracée en inclusion, le crochet du côté de $P\otimes S$, sans tête nette.}

\[
\alpha i = \beta i \;\overset{?}{=}\; (i\otimes\mathrm{id}_S)\circ\Delta_P
\]
\note{$\Delta_P$ est entouré, avec un point d'interrogation.}

$x\in C\otimes C'$, \quad \struck{$x$} $s\in S$, \quad $\alpha(x)=\beta(x)$ $\longrightarrow$ dans $\operatorname{Im}(x\otimes S \to C\otimes C'\otimes S)$ ; \quad $x\otimes S$.

\begin{tikzcd}[column sep=small, row sep=small, nodes={font=\scriptsize}]
 & & C'' \arrow[dl] \arrow[d] \arrow[dr] & \\
\cdots P_3 \arrow[r, hook] \arrow[d] & P_2 \arrow[r, hook] \arrow[d] & P_1 \arrow[r, hook] \arrow[dr] & C\otimes C' \arrow[d, bend right=20, "\alpha"'] \arrow[d, bend left=20, "\beta"] \\
P_2\otimes S \arrow[r, hook] & P_1\otimes S \arrow[rr, hook] & & C\otimes C'\otimes S
\end{tikzcd}

\note{la disposition est simplifiée : sur la page, $C''$ envoie trois flèches vers la chaîne $\cdots\subset P_3\subset P_2\subset P_1\subset C\otimes C'$, et chaque $P_{i+1}$ envoie une flèche oblique vers $P_i\otimes S$, $P_1$ vers $C\otimes C'\otimes S$.}

\[
P_\infty = \bigcap P_i
\]
$P_\infty$ \struck{$\otimes S$} $\longrightarrow C\otimes C'\otimes$ \uncertain{$C$}
\[
P_\infty \longrightarrow \bigcap P_i\otimes S = P_\infty\otimes S
\]

\[
\mathbf{P}
\]
\[
\text{\struck{$A\otimes$}}\quad S^h\otimes_h S^h = S, \qquad u\otimes 1 - 1\otimes u
\]
\note{la ligne $u\otimes 1-1\otimes u$ est suivie d'une croix (biffure ou renvoi) ; la lettre $h$ est lue sous réserve, ce pourrait être $k$.}

\[
C''\subset C\otimes_k C', \qquad \text{\struck{$M\otimes_k N$}}
\]
\[
u_C\otimes\mathrm{id}_{C'} = \mathrm{id}_C\otimes u_C
\]
\[
x\otimes 1 - 1\otimes x, \qquad x\otimes y - 1\otimes xy
\]
\[
V(C)\boxtimes_{V(S)} V(C')
\]

\uncertain{coproduits} dans \ill{}
\note{souligné.}

\[
C \longrightarrow C\boxtimes_S C
\]
\[
\text{\struck{$M\otimes M'\otimes S\otimes S$}}
\]
\[
M\otimes S \qquad M'\otimes S
\]
\[
V(C)\times V(C'') \longrightarrow V(C'')
\]
$V(S) \longrightarrow$ \ill{}

$\Longrightarrow$ $C''\longrightarrow C\boxtimes C'$ hom.\ de coalgèbres.

\page{284}
\begin{tikzcd}
S \arrow[r, "\Delta_S"] & S\otimes S & \\
C \arrow[u, "u"] \arrow[r, "\Delta_C"] & C\otimes C \arrow[u] \arrow[r, "{\mathrm{id}_C\otimes u}"] & C\otimes S
\end{tikzcd}

\note{les deux $S$ du haut sont écrits sur d'autres lettres ; l'étiquette $u$ de la flèche verticale gauche est reliée par un trait courbe.}

$V(S)$ \struck{Anneau} \uncertain{opère} sur $V(C)$

$V(C)$ structure d'algèbre sur $V(S)$ ?

$V(C)$ \struck{$\otimes$} \ill{} $V(C') \longrightarrow V(C'')$
bil.\ p.r.\ à $V(S)$

1°) bil.\ p.r.\ à \struck{\ill{}} \uncertain{$\odot$} d'un accouplement :
\[
C'' \longrightarrow C\otimes_k C'
\]

\begin{tikzcd}
V(C)\times V(S)\times V(C') \arrow[r, "{\mathrm{id}_{V(C)}\times\omega_{V(C')}}"] \arrow[d] & V(C)\times V(C') \arrow[d] \\
V(C'')\times V(S) \arrow[r] & V(C'')
\end{tikzcd}

\note{l'étiquette de la flèche verticale gauche, \struck{$\omega_{V(C)}\times\mathrm{id}_{V(S)}$}, est biffée.}

\[
V(C)\times\operatorname{Spec}(A) \longrightarrow V(C'') \qquad
\begin{array}{c} \text{\struck{$\uparrow$}} \\ C'' \end{array}
\qquad V(C)\longrightarrow V(C'')
\]
\[
\text{\struck{$V(C)\times V(S)\times\operatorname{Spec}(A)$}}
\]

\begin{tikzcd}
C \arrow[d] & C'' \arrow[l] \arrow[d] \\
C\otimes S & C''\otimes S \arrow[l]
\end{tikzcd}


\[
C\otimes A \longleftarrow C'' \qquad\qquad V(C)\times V(C') \longrightarrow V(C'')
\]
\[
\operatorname{Sym}(C)\otimes\operatorname{Sym}(C') \longleftarrow C''
\]

\begin{tikzcd}
C'' \arrow[r] \arrow[d] & C\otimes A \arrow[d] \\
C''\otimes S \arrow[r] & C\otimes S\otimes A
\end{tikzcd}

\struck{$C\otimes\operatorname{Sym}(C')\cap$} \ill{}
\[
\Delta_C\otimes\mathrm{id}_{C'} \simeq \mathrm{id}_C\otimes\Delta_{C'}
\]
\[
\simeq (\alpha\otimes\mathrm{id}_S)\circ\Delta_{C''}
\]

\begin{tikzcd}
C'' \arrow[r, "\alpha"] \arrow[d, "\Delta_{C''}"'] & C\otimes C' \arrow[d, "{\Delta_C\otimes\mathrm{id}_{C'}}"] \\
C''\otimes S \arrow[r, "{\alpha\otimes\mathrm{id}_S}"'] & C\otimes S\otimes C'
\end{tikzcd}

\note{une flèche courbe part de la droite du carré et revient sur $C\otimes S\otimes C'$ ; une autre relie la formule biffée $C\otimes\operatorname{Sym}(C')\cap\ldots$ à la ligne $\Delta_C\otimes\mathrm{id}_{C'}\simeq\ldots$.}

\section{Modules gradués : troncatures et foncteurs adjoints}

\page{286}
(1) $S$ anneau $\mathbb{Z}$-gradué \add{comm.} à \emph{degrés $\geq 0$}, $S=\sum_{n\geq 0} S_n$ \struck{$S_0=k$} ;

(2) $k = S_0$ ;

(3) $\operatorname{Modgr.d}(S) = \mathcal{M}$ ;

(4) $t$ foncteur translation
\[
\left\{\begin{aligned}
&t : \mathcal{M}\to\mathcal{M}\\
&(4.1)\quad t(M)_i = M_{i+1}
\end{aligned}\right.
\]
c'est un automorphisme de $\mathcal{M}$ ; $M[i] = t^i(M)$.

(4.2) \struck{$\mathcal{M}$} Pour tout entier $N$ on considère les deux sous-catégories pleines \add{strictes} $\mathcal{M}_N$, $\mathcal{M}^N$ de $\mathcal{M}$ définies par
\[
(5)\qquad
\begin{aligned}
\operatorname{Ob}\mathcal{M}_N &= \{P\in\operatorname{Ob}\mathcal{M} \mid P_i = 0 \ \ \forall i<N\}\\
\operatorname{Ob}\mathcal{M}^N &= \{P \mid P_i = 0 \ \ \forall i>N\}
\end{aligned}
\]

N.B. Les $\mathcal{M}_N$ ($\mathcal{M}^N$) décroissant (croissant) pour $N$ croissant. On a
\[
(6)\qquad t^i(\mathcal{M}_N) = \mathcal{M}_{N-i}, \qquad t^i(\mathcal{M}^N) = \mathcal{M}^{N-i}
\]
\struck{donc} En particulier \struck{est}
\[
(6.1)\qquad \mathcal{M}_N = t^{-N}\mathcal{M}_0, \qquad \mathcal{M}^N = t^{-N}\mathcal{M}^0 .
\]

\add{Donc} (La catégorie $\mathcal{M}$, munie déjà de l'automorphisme $t$, la donnée des familles $(\mathcal{M}_N)$ et $(\mathcal{M}^N)$ avec leurs inclusions déjà notées
\[
(7)\qquad \mathcal{M}_N\supset\mathcal{M}_{N'} \ \text{ et }\ \mathcal{M}^N\subset\mathcal{M}^{N'} \quad\text{si } N\leq N'
\]
\uncertain{équivaut} : \uncertain{à celle} du couple
\[
(8)\qquad \mathcal{M}_0,\ \mathcal{M}^0 \quad\text{avec}\quad t\mathcal{M}_0\subset\mathcal{M}_0,\quad t\mathcal{M}^0\supset\mathcal{M}^0 .
\]
\note{dans (8), le $\subset$ est repassé sur un autre signe.}
On se \uncertain{servira} d'ailleurs \uncertain{comme} $\mathcal{M}_N$, $\mathcal{M}^{N-1}$ \ill{}

\struck{\ill{}} \uncertain{se déterminent mutuellement} \ill{} $\mathcal{M}_0\simeq\mathcal{M}^0$ \uncertain{passant} \ill{} \uncertain{tous} \struck{\ill{}} $\mathcal{M}_N$, $\mathcal{M}^N$.
\note{ces deux dernières lignes sont écrites vite et très abrégées ; la première est en grande partie biffée.}

\page{288}
On considère les foncteurs d'inclusion
\[
(9)\qquad i_N : \mathcal{M}_N\hookrightarrow\mathcal{M}, \qquad i^N : \mathcal{M}^N\to\mathcal{M}
\]
plus généralement, pour $N\leq N'$
\[
(9.1)\qquad i_{N,N'} : \mathcal{M}_{N'}\hookrightarrow\mathcal{M}_N, \qquad i^{N',N} : \mathcal{M}^N\hookrightarrow\mathcal{M}^{N'}
\]
\struck{On se étudiera les}
\uncertain{d'où}
\[
(9.2)\qquad i_N\, i_{N,N'} = i_{N'}, \qquad i^{N'}\, i^{N',N} = i^N
\]
\struck{plus} et si $N\leq N'\leq N''$, d'où
\[
(9.3)\quad
\left\{\begin{aligned}
&\mathcal{M}_{N''}\subset\mathcal{M}_{N'}\subset\mathcal{M}_N \qquad \mathcal{M}^N\subset\mathcal{M}^{N'}\subset\mathcal{M}^{N''}\\
&i_{N,N''} = i_{N,N'}\, i_{N',N''} \qquad i^{N'',N} = i^{N'',N'}\, i^{N',N}
\end{aligned}\right.
\]

On va étudier la suite de foncteurs adjoints aux \uncertain{$i_N$, $i_{N,N'}$} d'une part, $i^N$, $i^{N,N'}$ d'autre part.

\emph{Proposition 1.} $i_N$ \add{$: \mathcal{M}_N\hookrightarrow\mathcal{M}$} s'insère dans une suite de 4 foncteurs adjoints
\[
(10)\qquad q_N \quad i_N \quad \tau_N \quad \varphi_N
\]
(chaque foncteur \uncertain{adjoint} pour $i_N$, $\tau_N$, $\varphi_N$ adjoint à droite de celui écrit à sa gauche) définis ainsi :

a) $\tau_N : \mathcal{M}\to\mathcal{M}_N$ (foncteur troncature à droite)
\[
\tau_N(P) \subset P \quad\text{(sous-objet de $P$)}
\]
\note{sous le $\subset$, un renvoi : « hom.\ d'adjonction » ; sous $\tau_N(P)$, un « $=$ » vertical renvoie à $i_N\tau_N(P)$.}
\[
(10.1)\qquad (\tau_N(P))_i = \begin{cases} P_i & \text{si } i\geq N\\ 0 & \text{si } i<N \end{cases}
\]

\page{290}
b) $\varphi_N : \mathcal{M}_N\to\mathcal{M}$
\[
(10.2)\quad
\left|\;\begin{aligned}
(\varphi_N(Q))_i &= \operatorname{Hom}^i_{\mathcal{M}}(\tau_{N-i}(S),\, Q)
\ \bigl(\longrightarrow \operatorname{Hom}^0_{\mathcal{M}}(S,Q) = Q_i\bigr)\\
&= \operatorname{Hom}^*_{\mathcal{M}}(\tau_{N-i}(S)[-i],\, Q)
\end{aligned}\right.
\]
\note{au-dessus du $Q$ de la première ligne, un « $=$ » vertical renvoie à $i_N Q$.}

\struck{Avec \ill{}} structure de module provenant (pour $k\in\mathbb{Z}$, $u\in S^*_k$) des applications
\[
\tau_{N-i}(S) \xleftarrow{\ \deg k\ } \tau_{N-i-k}(S)
\]
multiplication par $u$, en transposant ; on a un hom.\ (dans \struck{\ill{}} $\mathcal{M}$, \uncertain{décrit} dans (10.2))
\[
(10.2.1)\qquad \varphi_N(Q) \longrightarrow i_N(Q)
\]
d'où en appliquant $\tau_N$
\[
(10.2.2)\qquad \tau_N\varphi_N Q \xrightarrow[\text{adj.}]{\ \sim\ } Q
\]
qui est un isom.\ (c'est un hom.\ d'adjonction).

L'autre hom.\ d'adjonction
\struck{$\varphi_N\tau_N$…} $P \longrightarrow \varphi_N\tau_N P$
est déduit des \struck{\ill{}}
\[
(10.2.3)\quad
\left|\;\begin{aligned}
&\text{\struck{$\tau_N P \longrightarrow P$}} \quad \operatorname{Hom}^i(\tau_{N-i}S,\, P)\\
&P_i \longrightarrow (\varphi_N\tau_N P)_i \simeq \operatorname{Hom}^i(\tau_{N-i}S,\, \tau_N P)
\end{aligned}\right.
\]
\note{sur la page, $P_i$ envoie aussi une flèche oblique, marquée $\simeq$, vers $\operatorname{Hom}^i(S,P)$, d'où une flèche remonte vers $\operatorname{Hom}^i(\tau_{N-i}S,\tau_N P)$ ; dans $\tau_{N-i}S$ de la ligne du haut, l'indice est surchargé.}
\marginal{tout hom.\ de degré $i$ \struck{de $S$ dans $P$} de $\tau_{N-i}S$ dans $P$ est à valeurs dans $\tau_N P\subset P$}
\note{l'indice de $\tau_N P$ dans cette note est écrit sur un autre, biffé.}

c) (N.B. On n'utilisera pas $q_N$, on le \uncertain{met} par \uncertain{conscience}) :
\[
q_N(P) = P/\text{sous-module engendré par les } P_i \text{ pour } i<N
\]

\page{292}
\uncertain{Cor.}) Les foncteurs $i_N$, $\tau_N$ commutent à toutes les sortes de limites inductives et projectives, $q_N$ commute aux $\varinjlim$, $\varphi_N$ aux $\varprojlim$. Les foncteurs $i_N$, $\varphi_N$ sont pl.\ fid., les foncteurs $q_N$, $\tau_N$ sont des foncteurs localisation. \struck{\ill{}} Pour $\tau_N$ \uncertain{comme} il est exact, et commute aux $\varinjlim$ et $\varprojlim$) c'est la localisation par une sous-cat.\ abélienne épaisse, stable par $\varinjlim$ et $\varprojlim$ ; en fait c'est précisément la sous-cat.\ $\mathcal{M}^{N-1}$ (qui est ainsi décrite \struck{\ill{}} en termes de $\mathcal{M}_N$).
\note{tout le paragraphe est marqué d'un trait vertical en marge gauche ; dans « $q_N$, $\tau_N$ sont des foncteurs localisation », les deux symboles sont repassés sur d'autres.}

\page{294}
$C$ catégorie abélienne

$S$ anneau $\mathbb{Z}$-gradué comm.\ à \emph{degrés} $\geq 0$

$\mathcal{M}$ \add{$=\mathcal{M}(C,S)$} (catégorie des $C$-$S$-modules \emph{gradués}).

$\mathcal{M}_N$, $\mathcal{M}^N$ définis comme plus haut
\[
\mathcal{M}_N \xrightarrow{\ i_N\ } \mathcal{M}
\]
admet foncteur adjoint à droite
\[
\mathcal{M} \xrightarrow{\ \tau_N\ } \mathcal{M}_N, \qquad i_N\tau_N(P)\hookrightarrow P .
\]
\struck{qui admet}
Construisons adjoint à droite $\varphi_N$ de $\tau_N$ :
\[
\operatorname{Hom}_{\mathcal{M}_N}(\tau_N P,\, Q) \simeq \operatorname{Hom}_{\mathcal{M}}(P,\, \varphi_N Q)\ ?
\qquad (Q\in\mathcal{M}_N)
\]
Il faut \uncertain{bien} prouver que pour $Q$ fixé, le \uncertain{premier} membre, comme foncteur en $P$, soit représentable. Or il est clair que $\tau_N$ commute aux $\varinjlim$ quelconques, donc le foncteur \uncertain{envisagé} transforme les $\varinjlim$ en $\varprojlim$. Cela implique déjà ce qu'on veut lorsque $C$ \emph{ou} $C^\circ$ est une « cat.\ de \uncertain{Groth.} »…

Comment \add{alors} décrire les $(\varphi_N Q)_i\in\operatorname{Ob}C$ ? On a \struck{\ill{}} pour tout $M\in\operatorname{Ob}\mathcal{M}$
\[
M_i \simeq \mathbf{Hom}^*(S,\, M)_i
\]
\note{ici et dans la suite, $\mathbf{Hom}$ rend un Hom souligné.}
donc
\[
\begin{aligned}
(\varphi_N Q)_i &\simeq \mathbf{Hom}^*(S,\, \varphi_N Q)_i \simeq \mathbf{Hom}^*(S[-i],\, \varphi_N Q)_0\\
&\overset{\text{justifier}}{=} \mathbf{Hom}^*(\tau_N(S[-i]),\, Q)_0\\
&= \bigl(\mathbf{Hom}^*((\tau_{N-i}(S))[-i],\, Q)\bigr)_0
\end{aligned}
\]

\page{296}
\[
\begin{aligned}
&\text{\struck{$=\bigl(\mathbf{Hom}^*((\tau_N S)[-i],\, Q)\bigr)_0$}}\\
&\text{\struck{$=\mathbf{Hom}^*(\tau_{N-i}S$}}\\
&\simeq \bigl(\mathbf{Hom}^*(\tau_{N-i}(S),\, Q)\bigr)_i\\
&\simeq \mathbf{Hom}^i(\tau_{N-i}(S),\, Q)
\end{aligned}
\]
\note{les deux premières lignes sont barrées de traits obliques.}

Enfin si dans $C$ les $\varprojlim$ existent, alors $i_N$ admet un adjoint à gauche $q_N$.

En dualisant $C$, \uncertain{en} appliquant ce qui précède à $C^\circ$, en utilisant l'isom.
\[
\begin{array}{ccc}
\mathcal{M}(C,S)^\circ & \simeq & \mathcal{M}(C^\circ,S)\\
\cup & & \cup\\
\mathcal{M}^N(C,S)^\circ & \simeq & \mathcal{M}_{-N}(C^\circ,S)
\end{array}
\]
\note{l'exposant $N$ de $\mathcal{M}^N(C,S)^\circ$ est écrit sur un indice biffé ; l'indice du membre de droite porte un petit trait devant le $N$, lu comme un signe moins sous réserve.}
on voit que (si $C$ ou $C^\circ$ est une cat.\ de \uncertain{Groth.} par ex.) les foncteurs
\[
i^N : \mathcal{M}^N\longrightarrow\mathcal{M}
\]
s'insère dans \struck{admet} une suite de 4 \struck{d'adjoints} foncteurs adjoints
\[
\varphi^N \quad \tau^N \quad i^N \quad q^N
\]
qui se déduisent de façon \uncertain{évidente} \ill{}
\note{la phrase s'arrête là ; le bas de la page est vide.}

\page{298}
\note{Feuillet de calculs isolé, sans prose ; il reprend les troncatures dans un cadre où interviennent $S_k$ et un anneau de polynômes $k[T]$. Les notations n'y sont pas définies.}

\struck{$\delta\geq$} \struck{\ill{}}, \struck{$\beta\geq 0$,} $\alpha,\beta\geq 0$, $d\geq\beta-\alpha$
\note{le $0$ de $\alpha,\beta\geq 0$ est repassé sur une autre lettre, et le $d$ est écrit sur un autre symbole.}

\[
\operatorname{Hom}^d\bigl(\tau_{\alpha}M_{S_k},\, \tau_\beta N_{S_k}\bigr)
\xrightarrow{\ \sim\ } \operatorname{Hom}^d\bigl(\tau_{\alpha'}M_{S_k},\, \tau_{\beta'}N_{S_k}\bigr)
\]
\[
\left(\begin{aligned}
&\alpha'\geq\alpha\\
&\beta'\geq\beta\\
&\text{(biffé)}\\
&\beta'-\alpha'\geq d
\end{aligned}\right)
\]
\note{les indices de troncature de $M$ sont surchargés dans les deux membres ; la lecture $\alpha$, $\alpha'$ est tirée de la condition entre parenthèses et reste douteuse. Dans la parenthèse, une ligne biffée est illisible ; dans $\gamma\geq -d$, un symbole biffé précède $-d$ ; le « sinon » de la ligne suivante est une lecture douteuse.}
\[
\Big\downarrow{\scriptstyle 2}\qquad \gamma\geq{-d}\quad(\text{si } d\leq 0)
\]
\[
\operatorname{Hom}^d(\tau_\gamma M,\, N) \quad (\text{sinon}\ \operatorname{Hom}^d(M,\, \tau_\gamma N),\ \text{si } 0\leq\gamma\leq d,
\]
\[
\simeq\ \operatorname{Hom}^d\bigl(M_{S_k},\, \varphi_\beta\tau_\beta N_{S_k}\bigr) \quad (\text{si } d\leq 0)
\]
\[
\simeq\ \operatorname{Hom}^d\bigl(S_k,\, \varphi_0(\mathbf{Hom}_k(M,N)_{S_k})\bigr)
\]
\[
\simeq\ \bigl(\varphi_0(\mathbf{Hom}_k(M,N)_{S_k})\bigr)_d
\]
\note{dans l'avant-dernière ligne, $\mathbf{Hom}_k$ est écrit sur un premier Hom à indice biffé.}

$M_k$ $k$-module

$M_{S_k} = M\otimes_k k[T]$
\[
(\varphi_0 M_{S_k})_d = \operatorname{Hom}^{-d}(\tau_d S_k,\, M_{S_k})
\]
\[
\begin{array}{c} \cup \\ u \end{array}
\qquad
u(T^{(i)}) = \xi_i\, T^{(i-d)} \quad (i\geq d), \qquad \xi_i\in M
\]
\[
u(T^{(i)}T^{\ell}) = c_{i\ell}\,\xi_{i+\ell}\,T^{(i+\ell-d)}
\]
\note{sous $T^{(i)}T^{\ell}$, une accolade renvoie à $c_{i\ell}T^{(i+\ell)}$.}
\[
u(T^{(i)})\,T^{\ell} = \xi_i\, T^{(i-d)}\,T^{\ell} = c_{i-d,\ell}\,\xi_i\, T^{(i+\ell-d)}
\]
\[
i\geq d,\ \ell\geq 0 \;\Longrightarrow\; \boxed{\,c_{i,\ell}\,\xi_{i+\ell} = c_{i-d,\ell}\,\xi_i\,}
\]
\[
c_{d,\ell}\,\xi_{d+\ell} = \xi_d
\]
\note{la page s'arrête sur cette relation ; la suite du calcul n'est pas dans ce lot.}

\section{Notations ½ simpliciales. Constructions universelles}
\note{Ces deux lignes sont de sa main, seules sur un feuillet (p.~300) qui sert de couverture à la série suivante, hors de ce lot ; « ½ simpliciales » se lit « semi-simpliciales ».}

\end{document}
